% vntex sets the appendix names in Vietnamese; the report is in English.
\renewcommand{\appendixname}{Appendix}
\renewcommand{\appendixtocname}{Appendices}
\renewcommand{\appendixpagename}{Appendices}
\addto\captionsenglish{\renewcommand\chaptername{Appendix}}
\appendixtitleon
\appendixtitletocon
\begin{appendices}

%%%%%%%%%%%%%%%%%%%%%%%%%%%%%%%%%
\chapter{Detailed Review of LLM-Based Requirements Engineering Studies}
\label{apx:re-studies}

\begin{landscape}
\scriptsize
\setlength{\tabcolsep}{3pt}
\renewcommand{\arraystretch}{1.1}
\begin{longtable}{|p{2.0cm}|p{2.6cm}|p{2.6cm}|p{2.2cm}|p{4.0cm}|p{4.0cm}|p{4.4cm}|}
    \caption{Summary of AI techniques and evaluation results in requirements engineering.}
    \label{tab:apx-re-studies} \\
    \hline
    \textbf{Study} & \textbf{RE tasks} & \textbf{AI technologies} & \textbf{AI model} & \textbf{Dataset} & \textbf{Metric} & \textbf{Result} \\
    \hline
    \endfirsthead
    \hline
    \textbf{Study} & \textbf{RE tasks} & \textbf{AI technologies} & \textbf{AI model} & \textbf{Dataset} & \textbf{Metric} & \textbf{Result} \\
    \hline
    \endhead
    GPT-Powered~\cite{gorer2024gpt} &
    Requirements elicitation: generate interview scripts for practical RE training &
    RAG-like in-context learning; prompt chaining (outline-based sequential dialogue generation) &
    OpenAI GPT-4; custom GPT-4 agent &
    4 generated scenarios: meeting scheduler, social housing, digital health tracking, food delivery; a modified FIFA interview script as reference knowledge source &
    GRUEN for reference-free linguistic quality; expert evaluation of naturalness, coherence, completeness and RE interview practices (1--5) &
    GRUEN 0.79--0.81; naturalness 4/5, coherence 4/5, completeness 3/5; strong rapport and closing, weaker active listening and stakeholder identification \\
    \hline
    MARE~\cite{mare2024} &
    End-to-end RE: elicitation $\rightarrow$ modeling $\rightarrow$ verification $\rightarrow$ specification, performed iteratively and collaboratively &
    Multi-agent infrastructure; prompt engineering for specialized agents; shared workspace &
    GPT-3.5-turbo, text-davinci-002/003; MetaGPT &
    5 public cases (ATM, Cafeteria Ordering, Library, Assembly, Time Monitor); GoalModelDataset; 4 new cases (Smart Home, Temperature, Air Conditioner, Humidistat control) &
    Modeling: precision, recall, F1; specification: human evaluation of completeness, correctness, consistency (0--2) &
    F1 84.1\% (problem), 78.9\% (use case), 87.6\% (goal); SRS 0.98 completeness, 1.92 correctness, 1.98 consistency; multi-agent collaboration +1.1\% F1 over individual LLMs \\
    \hline
    LLMREI~\cite{korn2025llmrei} &
    Requirements elicitation &
    Autonomous LLM interview chatbot; zero-shot (LLMREI-short) and least-to-most prompting (LLMREI-long); fine-tuning explored, then abandoned &
    GPT-4o (2024-08-06); GPT-3.5 Turbo for fine-tuning only; WhisperX for transcription &
    33 simulated stakeholder interviews over 2 scenarios (hair/nail salon, ski resort); 50 interview transcripts from Ferrari et al.\ for fine-tuning &
    Likert ratings of 20 common interviewer mistakes; elicitation coverage against ground truth (not, partially, fully elicited); question adaptability categories &
    Mistakes comparable to trained human interviewers; up to 73.7\% of ground-truth requirements elicited; about half of the questions context-dependent; fine-tuned models underperformed \\
    \hline
    ROPE~\cite{ma2025rope} &
    Requirement elicitation and specification for LLM use; iterative requirement refinement from execution feedback &
    Interactive ROPE training suite; prompt optimization tool; execution engine &
    GPT-4o; o3-mini &
    6 ``LLM-hard'' replication tasks: 3 GPT applications and 3 coding games, in a counterbalanced pre/post-test &
    Requirement quality and output quality; learning gain (post-test minus pre-test) &
    +19.1\% learning gain vs.\ +0.7\% for conventional prompt engineering; requirement quality +25.4\%, output quality +12.7\% \\
    \hline
    XTRAREG~\cite{abualhaija2025llm} &
    Extraction of regulatory requirements &
    LangChain RAG pipeline; zero-shot and chain-of-thought prompts &
    GPT-4o; GPT-3.5-turbo; text-embedding-ada-002 &
    GDPR and additional legal sources; 108 reference requirements built manually &
    BLEU, ROUGE-L, METEOR, BERTScore; legal expert evaluation of correctness, groundedness, plausibility and coverage &
    GPT-4o clearly outperformed GPT-3.5, but coverage of the reference requirements was low and limited to the GDPR; performance drops without RAG; legal expert remains necessary \\
    \hline
    Hymel and Johnson~\cite{hymel2025analysis} &
    Early requirements elicitation &
    Single-shot prompt vs.\ one hour of expert work &
    ChatGPT 4.0 &
    50 business stakeholders, each submitting a software idea; one LLM and one human expert document per idea &
    Alignment score (1--10); completeness rating; source identification accuracy; t-tests &
    LLM documents scored higher and more consistently on alignment; no significant difference in completeness; both documents over-attributed to AI \\
    \hline
    QUARE~\cite{cheng2026quare} &
    Requirements quality analysis; cross-quality conflict negotiation; KAOS goal modeling; standards compliance &
    Multi-agent system; chain-of-thought; zero-shot; RAG; cosine similarity &
    gpt-4o-mini; bert-base-uncased &
    Autonomous driving (industrial); ATM and Library (MARE); RollCall and Bookkeeping (iReDev) &
    CHV, MDC, coverage uniformity, minimum axis coverage; conflict resolution rate, BERTScore F1; DAG topology, logical consistency, ISO/IEC/IEEE 29148 scores, compliance coverage &
    35.0 requirements per project; 98.2\% compliance coverage; 94.9\% semantic preservation; verifiability and feasibility 4.96/5 \\
    \hline
    VibeIntent~\cite{vibeintent2026} &
    Elicitation, analysis, backlog generation and validation &
    Multi-agent system; ReAct; persona and stakeholder simulation; human-in-the-loop; staged pipeline &
    GPT-5.4, GPT-5.4-mini &
    Cafe Booking, Textbook Exchange, Mood Tracking, Educational Feedback Platform, Fit Together; human and LLM-as-a-judge evaluation &
    Requirement count; MDC; Vendi Score; INVEST (1--5); ICC(2,k) &
    59 requirements; MDC 0.414; Vendi Score 14.39; INVEST 4.42/5; ICC 0.85 \\
    \hline
    Elicitron~\cite{ataei2025elicitron} &
    Requirements elicitation &
    LLM agents role-playing diverse users; serial vs.\ parallel agent generation with KMeans filtering; chain-of-thought &
    GPT-4-Turbo; text-embedding-ada-002 &
    Tent design case study with 20 human empathic lead user interviews as baseline; three experiments with 20 agents per condition &
    CHV, MDC, silhouette score; latent needs per condition (t-test); precision, recall, F1 of latent need detection &
    Serial generation yields the most diverse users; all conditions found more latent needs than the human baseline; criteria + CoT best for latent need detection \\
    \hline
    iReDev~\cite{jin2025iredev} &
    End-to-end RE &
    Multi-agent system; knowledge-driven agents; chain-of-thought; stakeholder simulation; event-driven communication; human-in-the-loop &
    GPT-4-turbo (2024-04-09) &
    10 real-world software systems (Bookkeeping, Random Roll Call, Shopping Site, Gomoku, Weather Forecast, Currency Converter, \ldots); manual ground truth and LLM-as-a-judge &
    Elicitation: CHV, MDC; modeling: precision, recall, F1, BLEU, BERTScore; specification: BLEU, BERTScore, G-Eval &
    CHV 0.47, MDC 0.62; model F1 0.389, BERTScore 0.593; SRS BERTScore 0.616, completeness 4.2, correctness 4.0, cohesiveness 4.1 \\
    \hline
    Almeida et al.~\cite{almeida2025elicitation} &
    Requirements analysis and classification &
    Prompt engineering; few-shot and persona prompting &
    GPT-4; DeepSeek-V3 &
    WEPGCOMP and Lugar; manual ground truth &
    Requirement count, precision, recall, F1 &
    GPT-4: P 75\%, R 61.76\%, F1 67.74\%; DeepSeek-V3: 60\%, 52.90\%, 56.30\%; both struggled with non-functional requirements \\
    \hline
    Binkhonain and Alfayez~\cite{binkhonain2025prompts} &
    Requirements classification &
    Zero-shot, few-shot, persona and chain-of-thought prompting; fine-tuning &
    Claude 3 Haiku, DeepSeek-V3, Gemini 2.0 Flash, GPT-4 Turbo, Llama 3.2 3B &
    PROMISE NFR; SecReq &
    Macro-F1; precision, recall, weighted and micro F1 &
    Few-shot better than zero-shot on all tasks; few-shot + persona (+ CoT) best; several configurations outperform fine-tuned BERT \\
    \hline
\end{longtable}
\end{landscape}

%%%%%%%%%%%%%%%%%%%%%%%%%%%%%%%%%
\chapter{Detailed Review of Persona and Stakeholder Simulation Studies}
\label{apx:persona-studies}

\begin{landscape}
\scriptsize
\setlength{\tabcolsep}{3pt}
\renewcommand{\arraystretch}{1.1}
\begin{longtable}{|p{2.2cm}|p{4.2cm}|p{2.8cm}|p{4.0cm}|p{5.0cm}|p{3.6cm}|}
    \caption{Summary of persona generation, coordination mechanisms, evaluation, and tested language models.}
    \label{tab:apx-persona-studies} \\
    \hline
    \textbf{Study} & \textbf{Persona / stakeholder generation} & \textbf{Number of agents} & \textbf{Coordination / interaction} & \textbf{Impact on performance / key findings} & \textbf{LLMs tested} \\
    \hline
    \endfirsthead
    \hline
    \textbf{Study} & \textbf{Persona / stakeholder generation} & \textbf{Number of agents} & \textbf{Coordination / interaction} & \textbf{Impact on performance / key findings} & \textbf{LLMs tested} \\
    \hline
    \endhead
    Persona Policies (PPol)~\cite{chopra2026beyond} &
    Evolutionary program search (OpenEvolve) optimizing a Python generator to discover behavioral axes (e.g.\ bursty, information gating) and expand them into task-preserving role-play policies &
    A population of $N$ persona policies per task ($N = 5, 8, 10$) &
    Multi-turn dialogue between persona-conditioned user simulator and a fixed assistant; Random Forest discriminator (human-likeness) and Chamfer distance (coverage) with reflection feedback &
    33--62\% absolute gains in fitness over default simulators; 80.4\% rated human in blinded evaluation; training on PPol traces raises out-of-distribution task success by 17\% &
    DeepSeek-V3.1, Qwen3-Next-80B, GPT-5.4-Mini; Gemma 4 31B; Gemini 3 Flash \\
    \hline
    Social Agents~\cite{bhattacharyya2026social} &
    Persona Agent Factory conditioning personas on demographic (age, gender, location) and psychographic (values, interests, lifestyle) attributes &
    $N = 10$ by default, up to 100 in scaling tests &
    Prompted independently without inter-agent communication to prevent groupthink; predictions aggregated by mean by a moderator &
    Better predictions on 11 tasks; +30.5\% on low-level and +9.9\% on high-level tasks; outperforms no-persona baselines and fine-tuned models &
    GPT-4o, LLaMA 3.x, Qwen3, Qwen2.5 VL \\
    \hline
    LLM Generated Persona~\cite{li2026persona} &
    Four tiers: meta (sampled census distributions), objective tabular, subjective tabular, descriptive (free-form narrative) &
    About 1{,}000 personas per state, 1{,}000{,}000 in total &
    Independent forced-choice survey simulation per persona; statistical population weighting; Wasserstein distance &
    More LLM-generated detail systematically increases bias and homogeneity, with a strong leftward political drift; overly positive profiles &
    Athene-70B, LLaMA-3.1-8B/70B, Mixtral-8x7B, Nemotron-70B, Qwen-2.5-72B, Yi-34B \\
    \hline
    PersonaGym~\cite{samuel2025personagym} &
    Persona agents initialized with persona-specific system prompts, evaluated in 150 persona-relevant environments &
    200 personas on 10{,}000 questions &
    Independent evaluation per question by an ensemble of two LLM evaluators guided by rubrics &
    PersonaScore; model size does not guarantee persona adherence; correlates with human judgment (Spearman 75.1\%) &
    GPT-3.5/4.1/4.5, LLaMA-2/3, Claude 3 Haiku, Claude 3.5 Sonnet, DeepSeek-V3 \\
    \hline
    MAJ-EVAL~\cite{chen2026majeval} &
    Two steps: extract evaluative dimensions and evidence from literature, then construct personas with demographics, evaluative dimension, specialty, psychological traits and social relationships &
    Panels organized by stakeholder groups (teachers, parents, children, clinicians, \ldots) &
    Independent evaluation, in-group free debate led by a moderator, out-group aggregator &
    Better alignment with human experts than ROUGE-L, BERTScore, G-Eval and ChatEval; in-group debate helps 15 of 20 stakeholder groups &
    Claude 3.7 Sonnet, Qwen-3-235B, GPT-4 \\
    \hline
\end{longtable}
\end{landscape}

%%%%%%%%%%%%%%%%%%%%%%%%%%%%%%%%%
\chapter{Agent Prompts}
\label{apx:prompts}
Every agent role has its own prompt, with a version number and the structure described in Section~\ref{sec:impl-prompts}. All prompts share a common section of settings, and every role also receives the output language. The stakeholder miner and the persona generator use the third version of their prompts, the other roles the second. Table~\ref{tab:apx-prompts} summarizes the inputs and outputs of each prompt; the prompt of the analyst is reproduced below as an example. Words in double braces, such as \{\{ transcripts \}\}, are replaced by the data of the run at each call.

\begin{table}[H]
    \centering
    \small
    \caption{Prompts of the agent roles.}
    \label{tab:apx-prompts}
    \begin{tabularx}{\linewidth}{|l|X|p{3.4cm}|}
        \hline
        \textbf{Role} & \textbf{Inputs} & \textbf{Output} \\
        \hline
        Planner & Product idea, current plan and the feedback of the user & Elicitation plan \\ \hline
        Corpus builder & Product idea, plan, number of documents to keep, search results & Source documents \\ \hline
        Stakeholder miner & Product idea, stakeholder groups already found, one source document & Stakeholder observations \\ \hline
        Persona generator & Product idea, plan, observations, aspects to cover, number of personas & Stakeholder groups, personas \\ \hline
        Interviewer & Plan, topics already covered, persona, transcript, turn, turn limit, rework directive & Question, or end of the interview with a summary \\ \hline
        Persona agent & Persona, what the persona knows, question, conversation & Answer \\ \hline
        Analyst & Plan topics, personas, transcripts, current analysis, rework directive & Analysis items \\ \hline
        Requirement engineer & Analyzed requirements, latent needs, assumptions, conflicts, current specification, directive & Specification requirements \\ \hline
        Backlog manager & Specification, current backlog, directive & Epics and user stories \\ \hline
        Validator & Specification, backlog, open issues, review round & Verdict, validation issues \\ \hline
        Intent router & Pending messages, recent history, project, specification, backlog & Target step and instruction \\ \hline
    \end{tabularx}
\end{table}

\section{Analyst Prompt}
\label{apx:analyst-prompt}
{\setstretch{1.0}
\lstinputlisting[style=prompt]{backmatter/prompts/analyst.v2.md}
}




\chapter{Project Management}
\label{apx:project-management}
\label{apx:timeline}
This appendix records how the specialized project was organized, so that its progress can be followed and the work of each member traced. It starts with an overview of the work areas and of the timeline, and then gives the details: the phases of the project and the log of the implementation. In every table, the column \emph{Assigned to} names the members responsible for each task.

The timeline and the log are reconstructed from the version history of the two repositories of the project: the repository of the system, which holds 92 commits between 28 September and 24 November 2026, and the repository of this report. The system was built with GitHub Flow: changes were made in pull requests, checked by the continuous integration of Section~\ref{sec:impl-testing}, and merged into the main branch as one commit each.

\section{Overview}
\label{apx:pm-overview}
The work of the project falls into six areas, listed in Table~\ref{tab:apx-areas}. Figure~\ref{fig:apx-gantt} shows when each area was active: the planning and the design came first; the web app was then built for about a month against simulated API responses, while the backend and the agents were developed; the backend joined the web app in November; and the report was written in the second half of November.

\begin{table}[H]
    \centering
    \caption{Work areas of the project.}
    \label{tab:apx-areas}
    \begin{tabularx}{\linewidth}{|p{3.4cm}|X|p{3cm}|}
        \hline
        \textbf{Area} & \textbf{Scope} & \textbf{Assigned to} \\
        \hline
    Planning and design & Goals, scope and key decisions; design documents of the architecture, the data, the API and the user interface & \\ \hline
    Web app (frontend) & Design system, screens and data access of the web app; simulated API; tests of the web app & \\ \hline
    Backend & API application, persistence, run execution, event log and streaming, exports, assistant & \\ \hline
    Agent workflow & Workflow definition, agent roles and prompts, search tool, integrity checks & \\ \hline
    Evaluation & Evaluation plan, scenarios and demonstration run & \\ \hline
    Report & Writing, diagrams and review of this report & \\ \hline
    \end{tabularx}
\end{table}

\begin{figure}[H]
    \centering
    \resizebox{\linewidth}{!}{% Timeline of the work areas of the specialized project (Appendix D), from the merged changes.
\begin{ganttchart}[
    time slot format=isodate,
    x unit=0.22cm, y unit chart=0.6cm, y unit title=0.55cm,
    title height=1, bar height=0.55,
    bar/.append style={fill=black!35, draw=black},
    milestone/.append style={fill=black},
    group/.append style={fill=black},
    bar label font=\small, group label font=\small\bfseries, title label font=\small,
    milestone label font=\small
  ]{2026-09-26}{2026-11-30}
  \gantttitlecalendar{month=shortname} \\
  \ganttbar{Planning}{2026-09-26}{2026-09-28} \\
  \ganttbar{Design documents}{2026-10-10}{2026-10-10}
  \ganttbar{}{2026-10-28}{2026-11-03}
  \ganttbar{}{2026-11-21}{2026-11-24} \\
  \ganttbar{Web app (frontend)}{2026-10-01}{2026-11-24} \\
  \ganttbar{Backend}{2026-11-10}{2026-11-24} \\
  \ganttbar{Agent workflow}{2026-11-13}{2026-11-24} \\
  \ganttbar{Report}{2026-09-26}{2026-09-26}
  \ganttbar{}{2026-11-16}{2026-11-27} \\
  \ganttmilestone{Backend joins the web app}{2026-11-10} \\
  \ganttmilestone{First version of the report}{2026-11-16}
\end{ganttchart}
}
    \caption{Timeline of the work areas of the specialized project, by date of the merged changes.}
    \label{fig:apx-gantt}
\end{figure}

\section{Timeline}
\label{apx:pm-timeline}
Table~\ref{tab:apx-timeline} details the phases of the project in each area.

\begingroup\setstretch{1.15}
\begin{longtable}{|p{0.15\linewidth}|p{0.12\linewidth}|p{0.45\linewidth}|p{0.16\linewidth}|}
    \caption{Timeline of the specialized project.}
    \label{tab:apx-timeline} \\
    \hline
    \textbf{Period (2026)} & \textbf{Area} & \textbf{Main results} & \textbf{Assigned to} \\
    \hline
    \endfirsthead
    \hline
    \textbf{Period (2026)} & \textbf{Area} & \textbf{Main results} & \textbf{Assigned to} \\
    \hline
    \endhead
    26--28 Sep & Setup & Repository of the report with its build and continuous integration & \\ \hline
    26--28 Sep & Setup & Planning document of the system: goals, scope, key decisions, risks and setup checklist; contributing guide & \\ \hline
    1--9 Oct & Web app & Skeleton of the web app with the component library and the design system & \\ \hline
    10 Oct & Design & Design documents of the architecture, the data shape, the API and the user interface & \\ \hline
    13 Oct -- 9 Nov & Web app & Screens and data access of the web app, built against simulated API responses (see Table~\ref{tab:apx-log}) & \\ \hline
    10 Nov & Backend & Backend behind the versioned API contract & \\ \hline
    13 Nov & Agents & Questions before planning, search tool of the corpus builder, retry of a failed stage & \\ \hline
    16--17 Nov & Report & Chapters laid out and written; UML diagrams; bibliography; status of the requirements & \\ \hline
    20--27 Nov & Report & Review of the background and related work; alignment with the system; evaluation plan and demonstration run & \\ \hline
    21--24 Nov & Web app & User and developer layouts; faces and acts of the agents & \\ \hline
    23--24 Nov & Backend & Assistant for questions about the application; fixes of authorization, evidence quotes and rework & \\ \hline
\end{longtable}
\endgroup

\section{Implementation Log}
\label{apx:pm-log}
Table~\ref{tab:apx-log} lists the main changes in the order in which they were merged, one task per row. Routine changes, such as refactoring, formatting and changes of the tooling, are left out.

\begingroup\setstretch{1.15}
\begin{longtable}{|p{0.12\linewidth}|p{0.12\linewidth}|p{0.48\linewidth}|p{0.16\linewidth}|}
    \caption{Implementation log of the project.}
    \label{tab:apx-log} \\
    \hline
    \textbf{Date (2026)} & \textbf{Area} & \textbf{Change} & \textbf{Assigned to} \\
    \hline
    \endfirsthead
    \hline
    \textbf{Date (2026)} & \textbf{Area} & \textbf{Change} & \textbf{Assigned to} \\
    \hline
    \endhead
    26 Sep & Report & Repository of the report with its local build and continuous integration & \\ \hline
    28 Sep & Planning & Planning document, README and contributing guide of the system & \\ \hline
    1 Oct & Web app & Skeleton of the web app with the component library, the chat shell and the test tooling & \\ \hline
    9 Oct & Web app & Design system and UI components & \\ \hline
    10 Oct & Design & Design documents of the architecture, the data shape, the API and the user interface & \\ \hline
    14 Oct & Web app & Workbench: activity bar, project workspace, tabs, top bar, status bar and command palette & \\ \hline
    15 Oct & Web app & Login, registration and sessions & \\ \hline
    15 Oct & Web app & Simulated API serving projects, threads and runs & \\ \hline
    18 Oct & Web app & Project list and project frame & \\ \hline
    21 Oct & Web app & Idempotency key on the requests that start work & \\ \hline
    27 Oct & Web app & Workflow definition, artifact kinds and content schemas & \\ \hline
    27 Oct & Web app & Reading and changing artifacts and items & \\ \hline
    27 Oct & Web app & Exports of a project & \\ \hline
    28 Oct & Web app & Trace links, traces and traceability views & \\ \hline
    29 Oct & Web app & Runs and run settings of a project & \\ \hline
    30 Oct & Web app & Review requests and their answers & \\ \hline
    31 Oct & Web app & Threads and messages & \\ \hline
    31 Oct & Web app & Run status and stop in the status bar & \\ \hline
    31 Oct & Web app & Model catalog and model settings of the user & \\ \hline
    31 Oct & Web app & Profile, password, sessions and settings of the user & \\ \hline
    1 Nov & Web app & Monthly usage and budget in the status bar & \\ \hline
    2 Nov & Web app & Outline and status of the documents; documents in the side panel and in tabs & \\ \hline
    3 Nov & Design & Changes of a document accepted only during a review; quotes of both items of a link & \\ \hline
    6 Nov & Web app & Editing, browsing and restoring records with their links and history & \\ \hline
    6 Nov & Web app & Documents drawn as layouts, conversations and profiles & \\ \hline
    7 Nov & Web app & Activity of a project: chat, panel and prompt & \\ \hline
    7 Nov & Web app & Sending messages, answering reviews and choosing run settings & \\ \hline
    9 Nov & Web app & Trace pages: linking, unlinking and opening items & \\ \hline
    10 Nov & Backend & Backend behind the versioned API contract & \\ \hline
    13 Nov & Agents & Questions asked before planning & \\ \hline
    13 Nov & Agents & Search tool of the corpus builder & \\ \hline
    13 Nov & Backend & Retry of a failed stage from its saved state & \\ \hline
    13 Nov & Web app & Phases of the work shown in time order & \\ \hline
    16 Nov & Report & Chapters laid out and written; diagrams redrawn in UML; bibliography; prompts in the appendix & \\ \hline
    20--22 Nov & Report & Review of the background and related work; alignment of the report with the system & \\ \hline
    21 Nov & Design & UI hints for kinds, personas and workflow acts & \\ \hline
    22 Nov & Web app & User and developer views, chosen by host & \\ \hline
    22 Nov & Web app & Faces of the agents and their acts & \\ \hline
    22 Nov & Web app & Chat with the reader beside it; pages of the user view & \\ \hline
    23 Nov & Web app & Two layouts of the desk; prompt with sources, dictation and a draft & \\ \hline
    23 Nov & Web app & Motions of the faces; acts in the chat; personas shown in blocks & \\ \hline
    23 Nov & Backend & Assistant that answers messages about the application & \\ \hline
    23 Nov & Backend & Fixes of authorization, evidence quotes, item conflicts and the assumption loop & \\ \hline
    24 Nov & Backend & Rework no longer copies the document and is told what failed & \\ \hline
    24 Nov & Design & Stakeholder types reduced to their word pattern & \\ \hline
    27 Nov & Report & Evaluation plan; agent descriptions updated; layout of the UML diagrams fixed & \\ \hline
\end{longtable}
\endgroup

\end{appendices}
